% Generated by tools/textbook/export_listings.js from tools/concepts/assessments.json: do not edit.
\begin{chapstatus}{Il filo del capitolo}
La macchina che avete programmato nella Lezione 02 è una macchina di von Neumann: programma e dati in un'unica memoria. RISC-V è l'ultimo passo di un lungo dibattito su quanto debbano essere semplici le istruzioni. Prossimo passo: come qualsiasi cosa diventa bit.
\par\medskip\textbf{Biglietto d'uscita} (da due a quattro minuti, da soli, su carta):
\begin{enumerate}[nosep,leftmargin=2em]
\item Spiegate la frase attribuita a Dijkstra, 'L'informatica non riguarda i computer più di quanto l'astronomia riguardi i telescopi', usando la definizione di informatica data nella lezione.
\item Dove è conservato il programma di una cartuccia mentre gira, e perché questo rende una macchina di von Neumann diversa da una macchina programmata ricollegando cavi?
\end{enumerate}
\end{chapstatus}
